\scanpage{31}

The condition a) already allows us to define the
relative $\pi_1$
\[
\pi_1(\sigma)
 = \pi_1(\Pi_{D^*}/\Pi_D),
\]
which is a local system of \unsure{profinite|profinite} abelian groups on $\Pi_D$
\inserted{\unsure{i.e.} a functor
$(\Pi_D)^{\mathrm{o}}\to\mathrm{Ens}$,}
and the last datum: \deleted{is a}

$4^\circ)$ A Kummerian isomorphism
\[
\deleted{\illegible{1}}\,k:\quad
\pi_1(\sigma)\simeq
\psi^*(\mathcal T)
\]

If we have two systems of this nature
\[
\Pi=(\Pi_U,\Pi_D,\Pi_{D^*},\varphi,\psi,\rho,\sigma,\alpha,k)
\]
and $\Pi'=(\Pi'_U,\Pi'_D,\ldots)$, a \uline{morphism}
from $\Pi'$ to $\Pi$ is a system of $3$ functors
\[
\begin{gathered}
f_U:\Pi'_U\to\Pi_U\\
f_D:\Pi'_D\to\Pi_D\\
\deleted{\illegible{1}}f_{D^*}:\Pi'_{D^*}\to\Pi_{D^*}
\end{gathered}
\]
and of \inserted{four}\deleted{\illegible{1}}
commutation isomorphisms
$\unsure{\alpha_{D^*D^*}|\alpha_{D^*D}},\alpha_{D^*U},\alpha_{Ue},\alpha_{De}$
for the four faces of the prism
\diagram{
\begin{tikzcd}[row sep=large,column sep=large]
 & \Pi'_{D^*}\arrow[rr,"\sigma'"]\arrow[dd,"\unsure{\rho'|\varphi'}"]
   \arrow[dl,"f_{D^*}"'] &&
   \Pi'_D\arrow[dd,"\psi'"]\arrow[dl,"f_D"']\\
 \Pi_{D^*}\arrow[rr,"\sigma",near start]\arrow[dd,"\rho"'] &&
 \Pi_D\arrow[dd,"\psi",near end]\\
 & \Pi'_U\arrow[rr,"\varphi'"]\arrow[dl,"f_U"'] &&
 \Pi_e\arrow[dl,"="']\\
 \Pi_U\arrow[rr,"\varphi"'] && \Pi_e
\end{tikzcd}
}
\deleted{\illegible{2}}

\scanpage{32}
\mspage{14}

satisfying an \inserted{equation} of compatibility with $\alpha,\alpha'$ which I do not write down
--- meaning that the \uline{two}
\deleted{\illegible{1}} isomorphisms \inserted{$u,v$} of functors
\diagram{[Two parallel arrows from $\Pi'_{D^*}$ to $\Pi_e$.
The upper arrow is labeled $\varphi\circ\rho\circ f_{D^*}$;
the lower arrow is labeled $\unsure{\mathrm{id}}\circ\psi'\circ\sigma'$.
Between the two, two downward arrows, labeled $u$ and $v$.]}
obtained respectively, $u$ by using successively
$\alpha,\alpha_{D^*D},\alpha_{De}$, \deleted{\unsure{and}} $v$ by
using successively $\alpha_{D^*U},\alpha_{Ue},\alpha'$, are equal.
[This compatibility could be expressed by interpreting the datum
of $\Pi$ as that of a fibered
category \inserted{$\Pi$} over the \deleted{\unsure{of diagram in the shape}}
\inserted{``category} square''
\[
Q:
\]
\diagram{[Diamond with vertices $D^*$ at the top, $\unsure{U}$ on the left,
$D$ on the right, $e$ at the bottom. The side $D^*\to U$ is labeled $\rho_0$;
the other three sides are labeled $\sigma_0$ between $D^*$ and $D$,
$\varphi_0$ between $U$ and $e$, $\psi_0$ between $D$ and $e$.
The arrowheads of these three lines are not clearly discernible.]}
where $\varphi_0\rho_0=\psi_0\sigma_0$.

with restriction to $\{e\}$ prescribed, and taking
\inserted{cartesian} functors between fibered categories\ldots].
Moreover, we require another compatibility, namely that

\scanpage{33}

\deleted{the homomorphism} the homomorphism of local systems of abelian groups on $\Pi'_D$
\[
(f_D)^*\bigl(\pi_1(\sigma)\bigr)
\longleftarrow \pi_1(\sigma')
\]
defined with the help of: $f_{D^*},f_D,\alpha_{D^*D}$
make the following diagram
of \deleted{\illegible{1}}isomorphisms \inserted{of}
\deleted{functors} local systems on $\Pi'_D$ commutative
\diagram{
\begin{tikzcd}[row sep=large,column sep=large]
(f_{D^*})^*\bigl(\pi_1(\sigma)\bigr)
 & \pi_1(\sigma')\arrow[l]\\
(f_{D^*})^*\bigl(\psi^*(\mathcal T)\bigr)
 \arrow[u,"(f_{D^*})^*(k)","\sim"']
 \arrow[d,"\mathrm{can}","\sim"']
 & \psi'^*(\mathcal T)
 \arrow[u,"k'","\sim"']\arrow[dl,"\sim"']\\
(\psi f_D)^*(\mathcal T) &
\end{tikzcd}
}
deduced from $\alpha_{De}:\psi'\simeq\psi\circ f_D$
\deleted{\illegible{1}}

\margin{\uline{no}, \unsure{that|this} only works for the
cases of étale morphisms, otherwise we must bring in
multiplication by
``degrees of ramification'' $d_{i'}$ $(i'\in\pi_0(\Pi'_{D^*}))$.}

For $\Pi,\unsure{\Pi^*|\Pi'}$ fixed, the systems
\[
(f,\alpha):=(f_U,f_D,f_{D^*},\alpha_{D^*D},\alpha_{D^*U},
\alpha_{U,e},\alpha_{D,e})
\]
above form a category in a natural \unsure{way}
--- in fact a groupoid --- by taking as morphisms
\inserted{$\mu$} from $(f',\alpha')$ to $(f,\alpha)$
the triples of morphisms
(profinite functors)
\[
f'_U\xrightarrow{\mu_U}\unsure{f_U},\qquad
f'_D\xrightarrow{\mu_D}f_D,\qquad
f'_{D^*}\xrightarrow{\mu_{D^*}}f_D
\]

\scanpage{34}
\mspage{15}

satisfying \deleted{some} four compatibility conditions with
$\alpha_{D^*D}$ and $\alpha'_{D^*D}$, with $\alpha_{D^*U}$ and
$\alpha'_{D^*U}$, with $\alpha_{Ue}$ and $\alpha'_{U,e}$, with
$\alpha_{D,e}$ and $\alpha'_{D,e}$ respectively.
(i.e. we are working in a full subcategory of the
\unsure{$2$}-category of the \uline{$\mathrm{Hom}$} between fibered
categories over $Q$, with fiber at $e$ fixed\ldots).
I have the impression that this groupoid
\uline{$\mathrm{Hom}$}$((f',\alpha'),(f,\alpha))$ is
always \uline{rigid}, i.e.
$\mathrm{Aut}(f,\alpha)$ is always reduced to the trivial group
--- I'm too lazy to check --- hence that if $(f',\alpha')$ and
$(f,\alpha)$ are isomorphic, the isomorphism in question is unique.
At any rate, we shall put

$\mathrm{Hom}((f',\alpha'),(f,\alpha))=\pi_0$
\uline{$\mathrm{Hom}$}$((f',\alpha'),(f,\alpha))$

Whence a \uline{category} of
\deleted{pairs \unsure{$(X,S)$}} systems
$(\Pi_U,\Pi_D,\Pi_{D^*},\varphi,\psi,\rho,\sigma,\alpha,
\deleted{\illegible{1}}k)$

\scanpage{35}

We have a functor from pairs $(X,S)$
(\unsure{i.e.|where} $X$ a scheme \inserted{locally} of finite type over $K$,
$S$ a \deleted{subset}\inserted{subscheme}
\inserted{closed \deleted{\illegible{1}}} \deleted{\unsure{discrete|discrete}} of $X$
\inserted{étale over $K$} (\unsure{such|such} that $\forall s\in S$, $X$
is smooth of relative dimension $1$ at $s$)
to this messy category $\mathcal B$.

\margin{\unsure{Morphisms|The morphisms} $(X',S')\to(X,S)$ are the
morphisms $f:X'\to X$ such that $f^{-1}(S)_{\text{reduced}}=S'$.}

\uline{Messy conjecture 7.1.}
When we restrict ourselves to the $(X,S)$ such that $X$ smooth projective of dimension
$\leq 1$, and which moreover are anabelian, then the
preceding functor is fully faithful.

\margin{\uline{N.B.} Faithfulness is easy\ldots}

Description of the messy category in terms of group theory:

Let $I=\pi_0(\Pi_{D^*})\simeq\pi_0(\Pi_D)$.
Let us choose an \unsure{element|element} \inserted{$D_i^*$} in each component of
$\Pi_{D^*}$, and let $D_i=\sigma(D_i^*)$.
For every $i$, let us choose an isomorphism
\[
\boxed{\psi(D_i)\xrightarrow[\simeq]{\unsure{\partial_{D_i}|\lambda_{D_i}}}
\mathrm{Spec}(\overline K)}.
\]
\deleted{\unsure{whence} \illegible{3}}
Up to replacing the object by an isomorphic ``subobject'',

\scanpage{36}
\mspage{16}

we may assume that $\Pi_{D^*}$ is the category sum of the categories
$[\mathcal E_i]$ defined by the
$\mathcal E_i=\mathrm{Aut}(D_i^*)$, $\Pi_D$
\deleted{is} \deleted{\illegible{1}}\inserted{the category}
sum of categories $[\Gamma_i]$ defined by the
$\Gamma_i=\mathrm{Aut}(D_i)$, the functor $\sigma$ being expressed
by a system of homomorphisms
\[
\sigma_i:\mathcal E_i\to\Gamma_i\qquad(i\in I)
\]
which are surjective with abelian kernels,
\inserted{the functor $\psi$ by a system of \uline{inclusions}
$\psi_i:\Gamma_i\subset\Gamma$}
and the datum of $k$ is in fact equivalent to isomorphisms
\[
k_i:\mathrm{Ker}\,\sigma_i\simeq T(\overline K)
\]
compatible with the operations of $\Gamma_i$ and $\Gamma$ on the two
\unsure{sides|sides} respectively and the inclusions $\psi_i$.

\margin{$^*$Again N.B. Two different bases \unsure{define}
different splittings!}

We can say that the datum of $(\sigma,\psi,k)$ is expressed
the datum of the system $(\mathcal E_i)_{i\in I}$, of a system
\unsure{of} \deleted{homomorphisms} exact sequences
\deleted{$p_i:\mathcal E_i\to\Gamma$}
\deleted{\illegible{5} and isomorphisms \illegible{1}}
\[
\boxed{1\to T(\overline K)\xrightarrow{k_i}\mathcal E_i
\xrightarrow{p_i}\Gamma}
\]
such that the $\Sigma_i
=\mathrm{Im}\,p_i$ are open, and that $k_i$ is compatible
with the operations of $\Sigma_i\simeq\mathrm{Coker}\,k_i$
(i.e. of $\mathcal E_i$ and $\Gamma$ on the two terms respectively).

\margin{\uline{N.B.} The extensions of the $\Gamma_i$ by
$T(\overline K)$ obtained from geometric situations
split --- \illegible{1} the choice of a uniformizer at $s_i$
defines a splitting, \unsure{and even} \illegible{1} the choice
of a basis of the tangent space at $s$\ldots}

\scanpage{37}

Assuming on the other hand (to simplify)
\deleted{\unsure{$X$}}\inserted{$\Pi_U$} connected,
\deleted{\unsure{i.e.} $U=X\setminus S$ connected, and}
and choosing an element $\widetilde U$ of $\Pi_U$ and an isomorphism
\deleted{between $\varphi(\widetilde U)$ and $\mathrm{Spec}\,K$}
\[
\varphi(\widetilde U)\xrightarrow{\unsure{\partial_U|\lambda_U}}
\mathrm{Spec}(\overline K)
\]
thus (up to replacing $\Pi_U$ by an equivalent groupoid)
we may assume $\Pi_U$ reduced to $\widetilde U$, and $\Pi_U$ is then given by
a group $\mathcal E$, and $\varphi$ by a group homomorphism
\[
\boxed{\mathcal E
\xrightarrow{\deleted{\illegible{1}}\,\unsure{\varphi_{\widetilde U,\partial_U}|\varphi_{\widetilde U,\lambda_U}}}
\Gamma}
\]
\inserted{(or $p$)}
whose image $\Sigma\subset\Gamma$ is thus an open subgroup
of $\Gamma$, and the kernel \unsure{will be denoted} $\unsure{\pi|\Pi}$
\[
\boxed{1\to\unsure{\pi|\Pi}\xrightarrow{\unsure{\kappa|\iota}}\mathcal E
\xrightarrow{\unsure{\varphi|p}}\Sigma_{\unsure{0}}\to1}
\]
\diagram{[Below $\Sigma_{\unsure{0}}$, a vertical inclusion sign toward $\Gamma$.]}

\margin{\unsure{Well noted} In geometric situations
\uline{N.B.} (This extension of group kernels splits, i.e.
$\exists$ an \unsure{open} subgroup $\Sigma_0$ \unsure{in} $\Gamma$,
\deleted{\unsure{of index} \illegible{1}} \inserted{and a} lifting
$\Sigma_0\to\mathcal E$\ldots}

[Having replaced the initial $\Pi_U$ by a smaller full subcategory,
we shall be obliged to modify $\rho$ up to isomorphism,
in practice by choosing for each $i\in I$ an isomorphism
\[
\rho(D_i)\xrightarrow{\mu_i}\widetilde U].
\]
\deleted{\illegible{1}}

\scanpage{38}
\mspage{17}

whereupon $\rho$ is made explicit simply by
group homomorphisms
\[
\boxed{\rho_i:\mathcal E_i\longrightarrow\mathcal E}
\]
It remains to make explicit the commutation isomorphism
\[
\alpha:\varphi\rho\to\psi\sigma
\]
which is defined by a system of elements
\[
\boxed{\gamma_i\in\Gamma\qquad(i\in I)}
\]
such that
\[
\unsure{(\alpha)|(*)}\qquad
\varphi\circ\rho_i=\mathrm{int}(\gamma_i)\circ p_i.
\]
which incidentally implies that $\rho_i$ \deleted{\illegible{1}}
maps $\mathrm{Ker}\,p_i$ into $\mathrm{Ker}\,p$, i.e. induces
a homomorphism of exact sequences
\diagram{
\begin{tikzcd}[row sep=large,column sep=large]
1\arrow[r] & T(\overline K)\arrow[r]\arrow[d] &
\mathcal E_i\arrow[r]\arrow[d,"\rho_i"] &
\deleted{\unsure{\mathcal E_i}}\inserted{\Sigma_i}
\arrow[r]\arrow[d,hook] & 1\\
1\arrow[r] & \unsure{\pi|\Pi}\arrow[r] & \mathcal E\arrow[r] &
\deleted{\unsure{\mathcal E_0}}\inserted{\Sigma_0}
\arrow[r] & 1
\end{tikzcd}
[To the right of $\Sigma_i$, as a superscript, $\subset\Gamma$; below $\Sigma_0$, an inclusion sign and $\Gamma$.]
}
with an induced injective homomorphism $\Gamma_i\to\Gamma_0$
--- and this being posited, the compatibility relation
\unsure{$(\alpha)$|$(*)$} becomes a relation on
group monomorphisms
\diagram{
\begin{tikzcd}
\Gamma_i\arrow[rr,hook]\arrow[dr] & & \Gamma\\
 & \Gamma_0\arrow[r,hook] & \Gamma
\end{tikzcd}
}

\scanpage{39}

\deleted{\uline{N.B.} When $I$ is empty i.e. $\Pi_{D^*},\Pi_D$
empty, then the ``messy'' systems \unsure{arising}
from geometric situations correspond to: $\unsure{\pi|\Pi}=\{1\}$, and
$\mathcal E$ reduces to a \unsure{subgroup}}

Thus, we have a relatively simple description of the
``reduced'' messy systems (i.e. where in the groupoids
$\Pi_U,\Pi_{D^*},\Pi_D$, each connected component has exactly
one object, and where \inserted{moreover} we force in some sense
$\unsure{\Pi_1(K)|\Pi_e(K)}$ to have only the object
$\mathrm{Spec}(\overline K)$.
But we are found again \unsure{on turning}
\unsure{for|by} the determination of the \uline{morphisms} of the systems
\[
\mathcal G=(\inserted{I,}G_i,p_i,k_i,\mathcal E,\varphi,\rho_i,
\gamma_i)
\quad\hbox{et}\quad
\mathcal G'=(I',G'_{i'},\ldots),
\]
\inserted{(or $p$?)}
\deleted{\illegible{2}} say in the \inserted{direction}
\[
f:\mathcal G'\to\mathcal G.
\]
We therefore need \inserted{(for $f_{D^*}$)} a map
\[
\boxed{\deleted{\illegible{1}}\,\unsure{\tau=\tau_f}:I'\to I},
\]
and $\forall\unsure{i'|i}$ a homomorphism
\[
\boxed{G'_{i'}\xrightarrow{f_{i'}}G_{\tau i'}},
\]
\deleted{inducing \unsure{by} \illegible{1}}
\inserted{(taking $f_D$ into account)}
inducing \inserted{\unsure{by passage to the}} quotient homomorphisms
\[
\Gamma'_{i'}\xrightarrow{g_{i'}}\Gamma_{\tau i'}
\]
\deleted{\unsure{but}}

\margin{The ``simplified messy objects'' are thus the
systems of homomorphisms of exact sequences
\diagram{
\begin{tikzcd}[row sep=large,column sep=large]
1\arrow[r] & T(\overline K)\arrow[r,"k_i"]
\arrow[d,"\unsure{\rho_i^0}"] &
\mathcal E_i\arrow[r,"p_i"]\arrow[d,"\rho_i"] &
\Gamma\arrow[d,"\mathrm{int}(\gamma_i)"]\\
1\arrow[r] & \unsure{\pi|\Pi}\arrow[r,"\unsure{\kappa}"] &
\mathcal E\arrow[r,"p"] & \Gamma
\end{tikzcd}
}}

\scanpage{40}
\mspage{18}

of open subgroups of $\Gamma$,
and the datum \inserted{$f_{D^*},f_D,$}
$\alpha_{D^*,D},\alpha_{D,e}$ is then equivalent
to the datum of elements $\alpha_i$
$(\unsure{i\in I|i'\in I'})$ of $\Gamma$,
\inserted{$\boxed{\alpha_{i'}\in\Gamma}
\quad(i'\in I')$}
such that
\deleted{$\forall i\in I$, $g\in\Gamma_i$, \illegible{2}}
\[
g_{i'}(\gamma')=\mathrm{int}(\alpha_{i'})(\gamma')
\qquad\forall i'\in I',\quad\gamma'\in\Gamma'_{i'}.
\]
The datum of $f_U$ is equivalent to the datum of a group homomorphism
\[
\boxed{f_{\mathcal E}:\mathcal E'\to\mathcal E}
\]
that of $\alpha_{D,e}$ is equivalent to the datum of
\[
\boxed{\alpha\in\Gamma}
\]
such that
\[
\unsure{(\alpha)|(*)}\qquad
\varphi f_{\mathcal E}=\mathrm{int}(\alpha)\varphi
\]
so that $f_{\mathcal E}$ induces a homomorphism of exact sequences
\diagram{
\begin{tikzcd}[row sep=large,column sep=large]
1\arrow[r] & \unsure{\pi|\Pi}'\arrow[r]\arrow[d] &
\mathcal E'\arrow[r]\arrow[d,"f_{\mathcal E}"] &
\Gamma'_0\arrow[r]\arrow[d,"\unsure{f_{\Gamma_0}|f_{\Sigma_0}}"] & 1\\
1\arrow[r] & \unsure{\pi|\Pi}\arrow[r] & \mathcal E\arrow[r] &
\Gamma_0\arrow[r] & 1
\end{tikzcd}
[Above $\Gamma'_0$, $\Gamma$ with a vertical inclusion
sign; below $\Gamma_0$, an inclusion sign and $\Gamma$.]
}
and with this, the condition called \unsure{$(\alpha)$|$(*)$}
\deleted{is}\inserted{becomes} a condition on
\inserted{inclusions of} subgroups of $\Gamma$:
\[
f_{\Gamma_0}(\gamma')=\mathrm{int}(\alpha)\gamma'
\quad\hbox{si}\quad\gamma'\in\Gamma'_0.
\]
We must still make explicit (in addition to the
$f_{D^*},f_D,f_U,\alpha_{D^*D},\alpha_{De},\alpha_{Ue}$
already made explicit) the